\documentclass[a4paper,11pt,reqno]{amsart}
\usepackage[utf8]{inputenc}
\usepackage{enumerate}
\usepackage{amsmath, amssymb,amsthm}
\usepackage{mathtools}
\usepackage{leftidx}
\usepackage{setspace}
\numberwithin{equation}{section}
%\usepackage{mathbbol}
\setcounter{tocdepth}{1} %2=show subsections, 1= show only sections
\usepackage[left=3cm, right=3cm]{geometry}
\usepackage{hyperref}
\usepackage{xcolor}
%\definecolor{name}{model}{color-spec}
\definecolor{my-black}{rgb}{0,0,0}
\definecolor{my-blue}{rgb}{0,0,0.8}
\definecolor{my-red}{rgb}{0.8,0,0} 
\definecolor{my-green}{rgb}{0,0.5,0}
\hypersetup{colorlinks,
    linkcolor	= {my-blue},
    citecolor	= {my-red},
    urlcolor		= {my-black}
}
\usepackage{graphicx}
\usepackage{tikz-cd}
\onehalfspacing
%Javier's suggestion to avoid overfull boxes
%\overfullrule=10cm better not to use because ArXiv might mess up things
%\usepackage{microtype}

\theoremstyle{plain} %italic

\newtheorem{lemma}{Lemma}[section]
\newtheorem{theorem}{Theorem}
\newtheorem{corollary}{Corollary}[section]
\newtheorem{proposition}{Proposition}[section]
\newtheorem{assumption}{Assumption}[section]
\newtheorem{problem}{Problem}[section]
\newtheorem{consequence}{Consequence}[section]
\newtheorem*{theorem*}{Theorem} %to avoid numbering


\theoremstyle{definition} %non-italic
\newtheorem{remark}{Remark}[section]

\newtheorem{definition}{Definition}[section]
\newtheorem{notation}{Notation}[section]
\newtheorem{conjecture}{Conjecture}[section]
\newtheorem{example}{Example}[section]
\newtheorem{claim}{Claim}[section]

\theoremstyle{remark}
\newtheorem*{remark-non}{Remark}

\DeclareMathOperator{\id}{id}
\DeclareMathOperator{\real}{Re}
\DeclareMathOperator{\imag}{Im}
\DeclareMathOperator{\image}{im}
\DeclareMathOperator{\supp}{supp}
\DeclareMathOperator{\vol}{vol}
\DeclareMathOperator{\inv}{inv}
\DeclareMathOperator{\Tr}{Tr}
\DeclareMathOperator{\ord}{ord}
\DeclareMathOperator{\Hom}{Hom}
\DeclareMathOperator{\Aut}{Aut}
\DeclareMathOperator{\End}{End}
\DeclareMathOperator{\Hol}{Hol}
\DeclareMathOperator{\GL}{GL}
\DeclareMathOperator{\SL}{SL}
\DeclareMathOperator{\PGL}{PGL}
\DeclareMathOperator{\PSL}{PSL}
\DeclareMathOperator{\Spur}{Sp}
\DeclareMathOperator{\PSO}{PSO}
\DeclareMathOperator{\SO}{SO}
\DeclareMathOperator{\Or}{O}
\DeclareMathOperator{\Res}{Res}
\DeclareMathOperator{\diam}{diam}
\DeclareMathOperator{\sgn}{sgn}
\DeclareMathOperator{\sinc}{sinc}
\DeclareMathOperator{\disc}{disc}


\newcommand{\R}{\mathbb{R}}
\renewcommand{\P}{\mathbb{P}}
\newcommand{\C}{\mathbb{C}}
\newcommand{\N}{\mathbb{N}}
\newcommand{\Q}{\mathbb{Q}}
\newcommand{\Z}{\mathbb{Z}}
\newcommand{\F}{\mathbb{F}}
\newcommand{\g}{\mathfrak{g}}
\renewcommand{\o}{\mathfrak{o}}
\renewcommand{\d}{\mathfrak{d}}
\renewcommand{\a}{\mathfrak{a}}
\renewcommand{\b}{\mathfrak{b}}
\newcommand{\p}{\mathfrak{p}}
\newcommand{\I}{\mathbf{I}}

\renewcommand{\sl}{\mathfrak{sl}}
\newcommand{\h}{\mathfrak{h}}
\renewcommand{\H}{\mathbb{H}}
\textwidth 6.5in
\oddsidemargin  -0.1in \evensidemargin -0.1in \parindent0.7em
\textheight 23cm



\title{SPR near Gaussian}



\begin{document}
\maketitle 

Let $F \in \mathcal{F}^2$ be such that $F^2 \in \mathcal{F}^2$. Then we have that 
\[
\|F^2 - F(0)^2\|_{\mathcal{F}^2}^2 \sim \int_{\C} |(F^2)'|^2 \, w(z) \, dz,
\]
where $w(z) = \frac{e^{-\pi|z|^2}}{1+\pi |z|^2}$. This follows by expanding out in power series. 

Then we get 
\begin{align*}
\int_{\C} |(F^2)'|^2 w(z) \, dz & \sim \int_{\C} |\nabla |F^2||^2 w(z) \, dz \cr 
 & = - \int_{\C} (|F|^2 - c) \Delta |F^2| w(z) \, dz - \int_{\C} (|F|^2 -c) \langle \nabla |F|^2, \nabla w\rangle \, dz \cr 
 & = - \int_{\C} (|F|^2-c)\Delta |F^2| w(z) \, dz + \frac{1}{2} \int_{\C} (|F|^2 - c)^2 \, \Delta w(z) \, dz. 
\end{align*}
The second term on the RHS is OK (as long as $c$ is the right thing), so we focus on the first one. We have 
\[
\int_{\C} (|F|^2 - c)\Delta|F^2| w(z) \, dz \le \left( \int_{\C} (|F|^2 - c)^2 \, d\mu\right)^{1/2} \left( \int_{\C} (\Delta|F^2|)^2 \frac{e^{-\pi|z|^2}}{P(|z|^2)} \, dz\right)^{1/2}, 
\]
where $P$ is a certain polynomial of degree $2$ with positive coefficients and $\mu$ is the Gaussian on $\C$. Let then $v$ be a real harmonic function. We have: 

\begin{align*}
J = \int_{\C} (\Delta|F|^2)^2 \frac{e^{-\pi|z|^2}}{P(|z|^2)} \, dz & = \int_{\C} \left(\Delta(|F|^2 - v)\right)^2 \, \frac{e^{-\pi|z|^2}}{P(|z|^2)} \, d z  = \int_{\C} (|F|^2-v)\Delta \left( \Delta(|F^2|) \frac{e^{-\pi|z|^2}}{P(|z|^2)}  \right) \, d z \cr 
    & = \int_{\C} (|F|^2 - v) \Delta^2(|F|^2) \frac{e^{-\pi|z|^2}}{P(|z|^2)} \, dz + \int_{\C} (|F|^2 -v) \nabla (\Delta |F|^2) \cdot \nabla \left( \frac{e^{-\pi|z|^2}}{P(|z|^2)} \right) \, dz \cr 
    & + \int_{\C} (|F|^2 -v) \Delta(|F|^2) \Delta\left( \frac{e^{-\pi|z|^2}}{P(|z|^2)} \right) \, dz = I_1 + I_2 + I_3.
\end{align*}
Note that 
$$\Delta(|F|^2) = 2 \left(|\nabla \text{Re}(F)|^2 + |\nabla \text{Im}(F)|^2 \right) = 4 |F'|^2,$$
and thus $\Delta^2(|F|^2) = 16 |F''|^2.$ Also, we have $|\nabla (\Delta (|F|^2))| \le C |F'||F''|, $ and $\left| \Delta \left( \frac{e^{-\pi|z|^2}}{P(|z|^2)}\right) \right| \le C \frac{e^{-\pi|z|^2}}{\sqrt{P(|z|^2)}}.$ Thus, 
\[
I_3 \le C \||F|^2 - v\|_{L^2_\mu} J^{1/2},
\]
\[
I_2 \le C \||F|^2 - v\|_{L^2_\mu} \left( \int_{\C} |F' \cdot F''|^2 \frac{e^{-\pi |z|^2}}{P(|z|^2)^{3/2}} \, dz \right)^{1/2},
\]
\[
I_1 \le C  \||F|^2 - v\|_{L^2_\mu} \left( \int_{\C} |F''|^4 \frac{e^{-\pi|z|^2}}{P(|z|^2)^2} \, dz \right)^{1/2}.
\]
We need now the following: 
\begin{lemma} The following bounds hold for any polynomial $F:$
\begin{align*}
\int_{\C} |F''|^4 \frac{e^{-\pi|z|^2}}{P(|z|^2)^2}\, dz & \le C \int_{\C} |F'(z)|^4 \frac{e^{-\pi|z|^2}}{P(|z|^2)} \, dz. 
\end{align*}
\end{lemma} 

\begin{proof} We prove only the second, since the first follows similarly. Let us write, using the Cauchy integral formula, 
\begin{align*}
\int_{|z|>20} |F''(z)|^4 \frac{e^{-\pi|z|^2}}{P(|z|^2)^2} \, dz & \le C \int_{|z| > 20} \left| \int_{|w| = 1/|z|} \frac{F'(w+z)}{w^2} \, \right|^4   \frac{e^{-\pi|z|^2}}{P(|z|^2)^2} \, dz \cr 
 & = C \int_{|z|> 20}  \left| \int_{0}^{2\pi} F'(e^{i\theta}/|z| \,+\,z) ie^{-2i\theta} \, \right|^4  |z|^4 \frac{e^{-\pi|z|^2}}{P(|z|^2)^2} \, dz \cr 
 & \le C \left( \int_0^{2\pi} \left( \int_{|z|>20} |F'(z + e^{i\theta}/|z|)|^4 \frac{e^{-\pi|z|^2}}{P(|z|^2)} \, dz \right)^{1/4} \, d\theta \right)^4 \cr 
 & \le C \int_{|w'|>10} |F'(w')|^4 \frac{e^{-\pi |w'|^2}}{P(|w'|^2)} \, dw'. 
\end{align*}
The contribution from the $|z|<20$ can be estimated in an even simpler way: the Cauchy integral formula can be used with the circle $\{|w| =1\}$ instead. 
\end{proof} 

Using the lemma and the computations before, we have 
\[
J \le C  \||F|^2 - v\|_{L^2_\mu}\cdot J^{1/2},
\]
as $J = \int_{\C} (\Delta|F^2|)^2 \, \frac{e^{-\pi|z|^2}}{1+\pi |z|^2} \, dz $ and $\Delta |F^2| \sim |F'|^2$. Hence, 
\[
J \lesssim \| |F|^2 - v\|_{L^2_\mu}^2.
\]
Selecting $v$ to be an appropriate constant finishes the proof. 

\end{document}

This implies that 
\[
J \le C \left( \||F|^2 - v\|_{L^2_\mu} + \||F|^2 - v\|_{L^2_\mu}^{1/2} \left( \||F|^2 - v\|_{L^2_\mu} + \|F\|_{L^4_\mu}^2\right)^{1/2} \right). 
\]
Notice now that, since $\Delta|F|^2 = 4 |F'|^2,$ then we have 
\[
J \ge \int_{\C} |F'(z)|^4 \frac{e^{-\pi|z|^2}}{P(|z|^2)} \, dz
\]
By Cauchy-Schwarz, thus,
\[
\int_{\C} |(F^2)'(z)|^2 \frac{e^{-\pi|z|^2}}{P(|z|^2)^{1/2}} dz \le \|F\|_{L^4_\mu}^2 J^{1/2}
\]
But now we redo the computations in \textsc{https://arxiv.org/pdf/2207.13606.pdf}: these show that, since $P^{-1/2}$ bounds the inverse of the first order Laguerre polynomial at the negative integers, then 
\[
\int_{\C} |H'(z)|^2 \frac{e^{-\pi|z|^2}}{P(|z|^2)^{1/2}} \, dz \ge c \int_{\C} |H'(z)|^2 \frac{e^{-\pi|z|^2}}{1+\pi|z|^2} \, dz \ge c' \min_{\alpha \in \C} \|H-\alpha\|_{L^2_\mu}^2. 
\]
Hence, putting it all together and taking $F = \frac{G}{\|G\|_{L^4_\mu}}$ and $v = \frac{\text{Re}(G)}{\|\text{Re}(G)\|_{L^2_\mu}},$ we have that, if 
\[
\left\| \frac{|G_k|^2}{\|G_k\|_{L^4_\mu}^2} - \frac{\text{Re}(G_k)}{\|\text{Re}(G_k)\|_{L^2_\mu}} \right\|_{L^2_\mu} \to 0,
\]
then 
\[
\min_{\alpha} \left\| \frac{G_k^2}{\|G_k\|_{L^4_\mu}^2} - \alpha \right\|_{L^2_\mu} \to 0.
\]
Since $\frac{G_k^2}{\|G_k\|_{L^4_\mu}^2}$ is a bounded sequence in the Fock space, $\alpha_k$ such that the minimum above is attained are also bounded, and thus, passing to a subsequence, can be taken to converge. Thus, 
\[
\frac{G_k^2}{\|G_k\|_{L^4_\mu}^2} \to \alpha_0 \text{ in } \mathcal{F}^2
\]
for some $\alpha_0 \in \C.$ But this implies uniform convergence on compact sets, and $\frac{G_k^2(0)}{\|G_k\|_{L^4_\mu}^2} = 0,$ which implies that $\alpha_0 = 0.$ This is a contradiction, since the sequence $\frac{G_k^2}{\|G_k\|_{L^4_\mu}^2} $ has norm $1$ in the Fock space. 

Thus, we conclude that \emph{no} sequence $G_k$ of polynomials can be such that 
\[
\left\| \frac{|G_k|^2}{\|G_k\|_{L^4_\mu}^2} - \frac{\text{Re}(G_k)}{\|\text{Re}(G_k)\|_{L^2_\mu}} \right\|_{L^2_\mu} \to 0,
\]
or, in other words, 
\[
\left\| \frac{|G|^2}{\|G\|_{L^4_\mu}^2} - \frac{\text{Re}(G)}{\|\text{Re}(G)\|_{L^2_\mu}} \right\|_{L^2_\mu} \ge c_0
\]
for some constant $c_0$ and all $G$ polynomials. Thus, by the improved version of Cauchy-Schwarz by Aldaz, the inequality 
\[
\int_{\C} \text{Re}(G) |G|^2 \, d\mu \le c \left( \int_{\C} \text{Re}(G)^2 \, d\mu \right)^{1/2} \left( \int_{\C} |G|^4 \, d\mu \right)^{1/2} 
\]
\emph{must} hold with some constant $c<1.$ \\

\noindent \textit{Remark.} Actually, no contradiction is needed in the end: since $G(0) = 0,$ we have 
\[ \min_{\alpha} \left\| \frac{G^2}{\|G\|_{L^4_\mu}^2} - \alpha \right\|_{L^2_\mu} \ge 1 \]
uniformly. Thus, It follows that there is a \emph{computable} $c_0 > 0$ such that 
\[
\left\| \frac{|G|^2}{\|G\|_{L^4_\mu}^2} - \frac{\text{Re}(G)}{\|\text{Re}(G)\|_{L^2_\mu}} \right\|_{L^2_\mu} \ge c_0
\]
for all polynomials $G.$ 

\newpage

Now, recall: we want to prove that 
\[
\int_{\C} \frac{f(z)^2}{(1+s \cdot f(z))^{3/2}} \, d\mu(z) \ge \frac{1}{2} \int_{\C} f(z)^2 d\mu(z),
\]
where $f(z) = 2r \text{Re}(G) + r^2|G|^2,$ whenever $s$ is sufficiently small (independently on $r,G$). 

Let then 
$$g(s) =\int_{\C} \frac{f(z)^2}{(1+s \cdot f(z))^{3/2}} \, d\mu(z).$$ 
Notice that 
$$g'(s) = -\frac{3}{2} \int_{\C} \frac{f(z)^3}{(1+s \cdot f(z))^{5/2}} \, d\mu(z),$$

and 

$$g''(s) = \frac{15}{4} \int_{\C} \frac{f(z)^4}{(1+s \cdot f(z))^{7/2}} \, d\mu(z).$$ 
First and foremost, notice that, since $f \ge -1,$ we have $g'' \ge 0$ for each $s \in (0,1).$ Thus, if we use Cauchy-Schwarz at the definition of $g',$ we will obtain 
\begin{align*} 
\left|\frac{2}{3} g'(s) \right| & = \left| \int_{\C} \frac{f(z)^3}{(1+s\cdot f(z))^{5/2}} \, d\mu(z) \right| \cr 
    & \le \left( \int_{\C} \frac{f(z)^2}{(1+s\cdot f(z))^{3/2}} \, d\mu(z)\right)^{1/2} \left( \int_{\C} \frac{f(z)^4}{(1+s\cdot f(z))^{7/2}} \, d\mu(z)\right)^{1/2} \cr 
    & = \left( g(s) \frac{4}{15} g''(s)\right)^{1/2}.
\end{align*} 
Hence, $g$ satisfies the following differential inequality: 
\[
\frac{5}{3} (g'(s))^2 \le g(s) g''(s).
\]
In other words, 
\[
\frac{2}{3} (g'(s))^2 \le \left( g(s)g''(s) - (g'(s))^2\right)
\]
\[
\iff \frac{2}{3} \left( \frac{g'(s)}{g(s)}\right)^2 \le \partial_s \left( \frac{g'(s)}{g(s)}\right). 
\]
Call $h(s) := g'(s)/g(s).$ Then these equations imply 
\[
-\frac{2}{3} \ge - \frac{h'(s)}{h^2} = \partial_s \left( \frac{1}{h(s)}\right). 
\]
Notice now that, since $g$ is \emph{convex} by the observations above, we may suppose that $g'(0) < 0,$ as otherwise $g$ is increasing and there is nothing to prove. In that case, let $I = (0,\alpha_0)$ be an interval where $g'(s) < 0, \forall s \in I.$ Then: 
\[
-\frac{2}{3} s \ge \frac{1}{h(s)} - \frac{1}{h(0)} = \frac{g(s)}{g'(s)} - \frac{g(0)}{g'(0)},
\]
and if $s \in I,$ we have $g'(s) g'(0) > 0,$ so 
\[
-\frac{2}{3} s g'(s)g'(0) \ge g(s) g'(0) - g'(s) g(0),
\]
\[
\iff g'(s) \left( g(0) - \frac{2}{3} s g'(0)\right) \ge g(s) g'(0). 
\]
As $g(0) > 0, g'(0) < 0,$ this implies, for $s \in (0,1),$
\[
\frac{g'(s)}{g(s)} = \partial_s \left( \log(g(s))\right) \ge \frac{g'(0)}{g(0) - \frac{2}{3} s g'(0)}. 
\]
Thus, 
\[
\log\left( \frac{g(t)}{g(0)} \right) \ge \int_0^t \frac{g'(0)}{g(0)-\frac{2}{3}g'(0) s} \, ds. 
\]
If $v = -\frac{2}{3} g'(0) s,$ then 
\[
\int_0^t \frac{g'(0)}{g(0)-\frac{2}{3}g'(0) s} \, ds = -\frac{3}{2} \int_0^{-\frac{2}{3} g'(0) t} \frac{dv}{g(0) + v} = -\frac{3}{2} \left( \log\left(g(0) - \frac{2}{3} g'(0) t\right) - \log(g(0))\right)\]
\[= -\frac{3}{2} \log \left( 1 - \frac{2}{3} t \frac{g'(0)}{g(0)}\right).
\]
Hence, 
\[
\frac{g(t)}{g(0)} \ge \left( 1 - \frac{2}{3} t \frac{g'(0)}{g(0)}\right)^{-3/2}. 
\]
Different idea: rewrite the differential inequality for $s \in I$ as 
\[
-\frac{5}{3} \frac{g'(s)}{g(s)} \le - \frac{g''(s)}{g'(s)},
\]
which is equivalent to 
\[
\frac{5}{3} \partial_s \left( \log(g(s))\right) \ge  \partial_s\left( \log(|g'(s)|)\right)
\]
\[
\iff \frac{2}{3} \partial_s \left( \log(g(s))\right) \ge \partial_s\left( \log(|g'(s)|/g(s))\right). 
\]
This implies 
\[
\frac{5}{3} \log(g(s)/g(0)) \ge \log(|g'(s)|/|g'(0)|),
\]
which is equivalent to 


\end{document} 